\chapter{System Model and Decision Policy Design}

% =====================================================================
% Capítulo portante de la contribución teórica. TODAS las variables,
% ecuaciones y criterios de decisión viven aquí; el Capítulo 4 únicamente
% los realiza. Regla de corte: si una ecuación aparece también en el
% Capítulo 4, sobra en uno de los dos.
% =====================================================================

\section{Scenario and System Model}
\emph{Definir formalmente el escenario: nodo móvil con doble interfaz, ruta con
paradas, infraestructura inalámbrica asociada a las paradas, cobertura celular
continua y servidor de recolección único. Introducir la notación que se usará en
todo el capítulo.}

\section{Assumptions and Design Constraints}
\emph{Enunciar explícitamente los supuestos: ruta conocida y repetitiva,
generación de datos por categorías con tasa conocida, memoria finita por
categoría, y disponibilidad de un historial de recorridos previos. Cada supuesto
debe justificarse por su plausibilidad operativa, no por conveniencia
matemática.}

\section{Characterisation of WiFi Communication Opportunities}
\emph{Definición conceptual clave de la tesis: la oportunidad de comunicación no
coincide con la parada física, sino con la ventana en que el nodo dispone
efectivamente de conectividad, que comienza antes de detenerse y persiste tras
reanudar la marcha. Formalizar la ventana y sus magnitudes.}

\section{Prediction of the Time to the Next Communication Opportunity}
\emph{Formular la estimación del tiempo hasta la próxima oportunidad y del tiempo
restante dentro de ella a partir de estadística histórica. Justificar aquí el
criterio de percentiles y su asimetría según el sentido del riesgo, y definir los
indicadores de fiabilidad asociados a la estimación.}

\section{Buffer Survival Condition and Minimum Offloading Volume}
\emph{Núcleo analítico de la política: formular, por categoría, la condición que
determina si la memoria disponible sobrevive hasta la siguiente evaluación y, en
caso negativo, el volumen mínimo que debe transmitirse por la interfaz costosa.
Introducir el horizonte de decisión y el factor de seguridad.}

\section{Deadline-Driven Rescue Condition}
\emph{Formular la segunda dimensión de la decisión: un dato debe transmitirse por
la interfaz costosa cuando la oportunidad prevista no llega dentro de su plazo de
validez, considerando un margen de guarda. Justificar por qué la condición de
supervivencia de memoria no cubre este caso.}

\section{Complete Formulation of the Offloading Policy}
\emph{Integrar ambas condiciones en una política única: precedencia entre ellas,
comportamiento en los casos límite y enunciado final del algoritmo en
pseudocódigo. Cerrar indicando qué propiedades se espera que exhiba, que son las
que el Capítulo 6 contrastará empíricamente.}
